\documentclass[12pt,a4paper]{article}
\usepackage[utf8]{inputenc}
\usepackage[T1]{fontenc}
\usepackage[spanish]{babel}
\usepackage{amsmath}
\usepackage{geometry}
\usepackage{array}
\usepackage{booktabs}
\usepackage{longtable}
\usepackage{float}
\usepackage{tikz}
\usetikzlibrary{arrows.meta,positioning,shapes.geometric,fit,backgrounds}
\geometry{left=3cm,right=2.5cm,top=2.5cm,bottom=2.5cm}
\setlength{\parindent}{1.25cm}
\setlength{\parskip}{0.35em}

\definecolor{azulDatos}{RGB}{219,235,255}
\definecolor{azulBorde}{RGB}{42,105,180}
\definecolor{verdeProc}{RGB}{222,246,229}
\definecolor{verdeBorde}{RGB}{40,135,88}
\definecolor{naranjaModelo}{RGB}{255,235,205}
\definecolor{naranjaBorde}{RGB}{211,120,32}
\definecolor{moradoMeta}{RGB}{237,226,255}
\definecolor{moradoBorde}{RGB}{116,80,170}
\definecolor{rojoSalida}{RGB}{255,226,226}
\definecolor{rojoBorde}{RGB}{190,65,65}
\definecolor{grisEval}{RGB}{238,241,245}
\definecolor{grisBorde}{RGB}{95,105,120}

\tikzset{
    bloque/.style={rectangle, rounded corners=5pt, draw=black, align=center, minimum height=1cm, text width=3.2cm},
    dato/.style={bloque, fill=azulDatos, draw=azulBorde, very thick},
    prep/.style={bloque, fill=verdeProc, draw=verdeBorde, very thick},
    modelo/.style={bloque, fill=naranjaModelo, draw=naranjaBorde, very thick},
    meta/.style={bloque, fill=moradoMeta, draw=moradoBorde, very thick},
    salida/.style={bloque, fill=rojoSalida, draw=rojoBorde, very thick},
    eval/.style={bloque, fill=grisEval, draw=grisBorde, very thick},
    flecha/.style={-Latex, very thick, draw=grisBorde},
    flechaAzul/.style={-Latex, very thick, draw=azulBorde},
    flechaVerde/.style={-Latex, very thick, draw=verdeBorde},
    flechaNaranja/.style={-Latex, very thick, draw=naranjaBorde},
    flechaMorada/.style={-Latex, very thick, draw=moradoBorde},
    flechaRoja/.style={-Latex, very thick, draw=rojoBorde}
}

\begin{document}

\setcounter{section}{3}
\section{Materiales y métodos}

\subsection{Tipo y diseño de la investigación}

\subsubsection{Tipo de investigación}

\subsubsection{Diseño de investigación}

El diseño de investigación fue experimental puro mediante simulación en entorno controlado. La experimentación se desarrolló sobre el dataset CICIoT2023, utilizando particiones independientes de entrenamiento, validación y prueba. El grupo experimental estuvo constituido por el modelo híbrido basado en Stacking, cuyos clasificadores de Nivel 0 fueron XGBoost y LightGBM, y cuyo meta-modelo de Nivel 1 fue un Árbol de Decisión alimentado por probabilidades de clase y metadatos de incertidumbre obtenidos mediante entropía de Shannon.

El grupo de control estuvo conformado por modelos base no hibridados, específicamente Random Forest y XGBoost. Estos modelos permitieron establecer una línea base de comparación para determinar si la arquitectura híbrida propuesta incrementa la detección de ataques minoritarios y mantiene una eficiencia computacional viable.

La representación del diseño experimental fue la siguiente:

\[
GE: X \rightarrow O_1
\]

\[
GC: - \rightarrow O_2
\]

Donde $GE$ representa el grupo experimental, $X$ representa la aplicación del modelo híbrido basado en Stacking con Borderline-SMOTE, $O_1$ representa las métricas obtenidas por el modelo propuesto, $GC$ representa el grupo de control y $O_2$ representa las métricas obtenidas por los modelos base.

\begin{figure}[H]
\centering
\resizebox{0.98\textwidth}{!}{%
\begin{tikzpicture}[node distance=1.0cm and 1.0cm]
\node[dato] (data) {CICIoT2023\\registros de flujo\\47 variables};
\node[prep, right=of data] (prep) {Limpieza\\codificación\\partición};
\node[prep, right=of prep] (smote) {Borderline-\\SMOTE\\solo train};

\node[modelo, below left=1.35cm and 0.25cm of smote] (ge) {Grupo experimental\\Stacking\\XGBoost + LightGBM};
\node[modelo, below right=1.35cm and 0.25cm of smote] (gc) {Grupo control\\Random Forest\\XGBoost};
\node[meta, below=2.4cm of smote] (val) {Evaluación común\\Test independiente};
\node[eval, below=of val] (metrics) {Recall\\F1-Score\\Accuracy\\latencia\\memoria};

\draw[flechaAzul] (data) -- (prep);
\draw[flechaVerde] (prep) -- (smote);
\draw[flechaNaranja] (smote) -- (ge);
\draw[flechaNaranja] (prep) |- (gc);
\draw[flechaMorada] (ge) -- (val);
\draw[flechaMorada] (gc) -- (val);
\draw[flechaRoja] (val) -- (metrics);
\end{tikzpicture}
}
\caption{Diseño experimental para la comparación entre el modelo híbrido y los modelos base.}
\end{figure}

\subsection{Población y muestra}

\subsubsection{Población}

La población estuvo constituida por los registros de tráfico de red IoT contenidos en el dataset CICIoT2023, el cual integra flujos benignos y flujos asociados a ataques informáticos generados en un entorno de prueba con dispositivos IoT. Cada registro representó una unidad de análisis compuesta por variables numéricas de flujo de red y una etiqueta de clase que identifica el comportamiento benigno o el tipo de ataque correspondiente.

En términos metodológicos, la población de interés correspondió a los registros válidos de tráfico normal y malicioso que permiten evaluar sistemas de detección de intrusiones basados en aprendizaje automático. Por tanto, no se trabajó con usuarios humanos, sino con datos secundarios etiquetados y estandarizados para investigación experimental en ciberseguridad.

\subsubsection{Muestra}

La muestra fue censal respecto al conjunto de datos procesado para el experimento, debido a que se emplearon todos los registros válidos disponibles en las particiones definidas para entrenamiento, validación y prueba. La distribución utilizada fue la siguiente:

\begin{table}[H]
\centering
\caption{Distribución de los datos empleados en la investigación}
\begin{tabular}{lrr}
\toprule
\textbf{Subconjunto} & \textbf{Registros} & \textbf{Variables} \\
\midrule
Entrenamiento & 5\,491\,971 & 47 \\
Validación & 1\,176\,851 & 47 \\
Prueba & 1\,176\,851 & 47 \\
\midrule
Total & 7\,845\,673 & 47 \\
\bottomrule
\end{tabular}
\end{table}

El tipo de muestreo fue censal no probabilístico, dado que la selección de registros no dependió de un sorteo de unidades, sino del uso completo de las particiones disponibles luego del preprocesamiento. Esta decisión fue pertinente porque el objetivo de la investigación fue evaluar el comportamiento del modelo sobre un volumen masivo de tráfico IoT y no inferir opiniones o características de una población humana.

\subsubsection{Criterios de inclusión y exclusión}

Se consideraron como criterios de inclusión los registros que presentaron etiqueta de clase válida, valores compatibles con el esquema de 47 variables del dataset y pertenencia a las particiones de entrenamiento, validación o prueba. También se incluyeron registros benignos y maliciosos, debido a que la detección de ataques minoritarios requiere contrastar las clases de ataque frente al comportamiento normal y frente a otras clases de amenaza.

Se excluyeron registros duplicados no informativos, registros con etiquetas ausentes, valores no numéricos incompatibles con los algoritmos seleccionados, instancias con valores infinitos o no representables y columnas identificadoras que no aportaron información predictiva al modelo. La exclusión tuvo como finalidad evitar ruido, fuga de información y sesgos artificiales durante el entrenamiento.

\subsection{Unidad de análisis}

La unidad de análisis fue el flujo de tráfico de red IoT representado por un registro del dataset CICIoT2023. Cada flujo estuvo compuesto por 47 variables predictoras y una etiqueta de clase. En el experimento, cada unidad de análisis fue procesada por los clasificadores para determinar si correspondía a tráfico benigno o a una categoría de ataque.

La unidad de observación coincidió con la unidad de análisis, puesto que las mediciones se realizaron directamente sobre cada registro de tráfico. A partir de las predicciones obtenidas para cada flujo se construyeron matrices de confusión, métricas por clase y métricas globales, priorizando el comportamiento de las clases minoritarias.

\subsection{Materiales}

\subsubsection{Dataset}

El material empírico principal fue el dataset CICIoT2023, seleccionado por su pertinencia para la detección de intrusiones en redes IoT. Este conjunto contiene tráfico benigno y tráfico malicioso correspondiente a escenarios modernos de ataque. Su estructura etiquetada permitió aplicar aprendizaje supervisado, mientras que sus características de flujo hicieron posible entrenar modelos tabulares como Random Forest, XGBoost, LightGBM y Stacking.

\subsubsection{Hardware}

El desarrollo experimental se planteó para ejecutarse en una laptop de alto rendimiento con procesador Intel Core i7, 32 GB de memoria RAM y tarjeta gráfica NVIDIA RTX 4060 con soporte CUDA. Aunque los modelos principales se basan en árboles y pueden ejecutarse en CPU, la disponibilidad de GPU y memoria suficiente favorece el procesamiento por lotes, la carga de grandes volúmenes de datos y la reducción del tiempo de experimentación.

\subsubsection{Software}

El entorno de software estuvo conformado por Python 3.11 y librerías especializadas para análisis de datos, aprendizaje automático, balanceo de clases y evaluación estadística. Las herramientas consideradas fueron Pandas y NumPy para manipulación de datos; Scikit-Learn para particionamiento, métricas, Random Forest, Árbol de Decisión y utilidades de validación; XGBoost para el clasificador base de aumento de gradiente extremo; LightGBM para el clasificador base optimizado por crecimiento de hojas; Imbalanced-learn para Borderline-SMOTE; SciPy para la prueba de Wilcoxon; y Matplotlib o Seaborn para la representación gráfica de resultados.

\begin{table}[H]
\centering
\caption{Materiales e instrumentos tecnológicos}
\begin{tabular}{p{4cm}p{8.5cm}}
\toprule
\textbf{Componente} & \textbf{Especificación} \\
\midrule
Dataset & CICIoT2023 con registros benignos y de ataque en redes IoT. \\
Lenguaje & Python 3.11. \\
Librerías principales & Pandas, NumPy, Scikit-Learn, XGBoost, LightGBM, Imbalanced-learn, SciPy, Matplotlib y Seaborn. \\
Hardware & Intel Core i7, 32 GB RAM, NVIDIA RTX 4060 con soporte CUDA. \\
Indicadores & Recall, F1-Score, Accuracy, tiempo de entrenamiento, latencia de inferencia y memoria RAM. \\
\bottomrule
\end{tabular}
\end{table}

\subsection{Método de recolección de datos}

El método de recolección fue documental y experimental. Fue documental porque se utilizó un dataset público y estandarizado, previamente etiquetado para investigación en seguridad IoT. Fue experimental porque los datos fueron sometidos a un proceso controlado de entrenamiento, validación y prueba con algoritmos de aprendizaje automático, obteniendo métricas cuantitativas de desempeño.

La técnica de recolección consistió en la descarga, organización, lectura y verificación del dataset CICIoT2023. El instrumento de recolección fue una ficha técnica de datos, en la cual se registraron la procedencia del conjunto de datos, número de registros, número de variables, clases disponibles, distribución de particiones, criterios de limpieza y transformaciones aplicadas.

\subsection{Procedimiento metodológico}

El procedimiento siguió la metodología CRISP-DM, adaptada al desarrollo de modelos de detección de intrusiones en redes IoT. Esta metodología fue seleccionada porque permite organizar el ciclo experimental en fases ordenadas, desde la comprensión del problema hasta la evaluación del modelo.

\subsubsection{Comprensión del problema}

En esta fase se delimitó el problema de investigación: la baja eficacia de los modelos de detección frente a ataques minoritarios en redes IoT. Se identificó que el desbalance de clases puede provocar altos falsos negativos, por lo que se definió como prioridad maximizar el Recall y el F1-Score de las clases minoritarias sin descuidar la eficiencia computacional.

\subsubsection{Comprensión de los datos}

Se revisó la estructura del CICIoT2023, identificando número de registros, variables predictoras, etiquetas de clase, distribución de ataques y presencia de desbalance. También se verificó la separación en entrenamiento, validación y prueba, con el propósito de evitar fuga de información durante el proceso experimental.

\subsubsection{Preparación de datos}

La preparación incluyó limpieza de registros, tratamiento de valores inválidos, codificación de etiquetas, separación de variables predictoras y variable objetivo, y verificación de consistencia dimensional. Borderline-SMOTE se aplicó únicamente al conjunto de entrenamiento, debido a que aplicar sobremuestreo sobre validación o prueba alteraría la distribución real de evaluación y produciría una medición artificial del rendimiento.

\subsubsection{Modelado}

En la fase de modelado se entrenaron los modelos base de control y el modelo experimental. Los modelos de control fueron Random Forest y XGBoost sin arquitectura de apilamiento. El modelo experimental estuvo compuesto por XGBoost y LightGBM como clasificadores base de Nivel 0. Posteriormente, sus probabilidades de clase fueron utilizadas para construir metacaracterísticas, a las que se añadió la entropía de Shannon como medida de incertidumbre. Finalmente, un Árbol de Decisión actuó como meta-modelo de Nivel 1 para emitir la predicción final.

\subsubsection{Evaluación}

La evaluación se realizó sobre el conjunto de prueba, que no participó en el entrenamiento ni en el ajuste del meta-modelo. Se calcularon métricas de clasificación y métricas computacionales. La comparación entre el modelo híbrido y los modelos base permitió determinar si la propuesta mejora la detección de ataques minoritarios y si mantiene tiempos de inferencia viables.

\begin{figure}[H]
\centering
\resizebox{0.98\textwidth}{!}{%
\begin{tikzpicture}[node distance=1.0cm and 0.95cm]
\node[dato] (p1) {1. Comprensión\\del problema};
\node[dato, right=of p1] (p2) {2. Comprensión\\de datos};
\node[prep, right=of p2] (p3) {3. Preparación\\de datos};
\node[modelo, below=of p3] (p4) {4. Modelado\\GE y GC};
\node[eval, left=of p4] (p5) {5. Evaluación\\métricas};
\node[salida, left=of p5] (p6) {6. Resultados\\y contraste};
\draw[flechaAzul] (p1) -- (p2);
\draw[flechaVerde] (p2) -- (p3);
\draw[flechaNaranja] (p3) -- (p4);
\draw[flechaMorada] (p4) -- (p5);
\draw[flechaRoja] (p5) -- (p6);
\draw[flecha] (p5.north) |- node[pos=0.25, above]{ajuste si corresponde} (p3.south);
\end{tikzpicture}
}
\caption{Procedimiento metodológico basado en CRISP-DM para la investigación.}
\end{figure}

\subsection{Arquitectura experimental del modelo propuesto}

La arquitectura experimental se estructuró en dos niveles. En el Nivel 0 se entrenaron XGBoost y LightGBM sobre los datos de entrenamiento previamente procesados y balanceados. Ambos algoritmos generaron probabilidades de clase para cada registro. Estas probabilidades no fueron interpretadas como salida final, sino como información intermedia para el meta-modelo.

En el Nivel 1 se construyó una matriz de metacaracterísticas compuesta por las probabilidades de clase de XGBoost, las probabilidades de clase de LightGBM y la entropía de Shannon calculada sobre dichas distribuciones probabilísticas. El Árbol de Decisión recibió estas metacaracterísticas y aprendió reglas de combinación para clasificar el tráfico como benigno o como ataque.

La arquitectura se expresa de la siguiente forma:

\[
M_{hibrido}=DT(P_{XGB},P_{LGBM},H)
\]

Donde $M_{hibrido}$ representa el modelo final, $DT$ representa el Árbol de Decisión de Nivel 1, $P_{XGB}$ corresponde a las probabilidades generadas por XGBoost, $P_{LGBM}$ corresponde a las probabilidades generadas por LightGBM y $H$ representa la entropía de Shannon como metadato de incertidumbre.

\subsection{Técnicas de procesamiento y análisis de datos}

El procesamiento de datos se realizó por lotes debido al volumen de registros. Esta estrategia permitió cargar y transformar los datos de manera controlada, reduciendo el riesgo de saturación de memoria. Se aplicaron técnicas de limpieza, verificación de tipos de datos, transformación de etiquetas y balanceo de clases mediante Borderline-SMOTE.

El análisis de datos incluyó estadística descriptiva de la distribución de clases antes y después del rebalanceo, generación de matrices de confusión, cálculo de métricas de clasificación y comparación de rendimiento entre el grupo experimental y el grupo de control. Asimismo, se registraron los tiempos de entrenamiento e inferencia para evaluar la viabilidad operativa del modelo propuesto.

\subsection{Métricas de evaluación}

Las métricas principales fueron Recall, F1-Score y Accuracy. El Recall fue priorizado porque mide la proporción de ataques reales correctamente detectados, lo cual es crítico cuando se evalúan clases minoritarias. El F1-Score permitió analizar el equilibrio entre Precision y Recall, evitando seleccionar modelos que detecten ataques a costa de un número excesivo de falsas alarmas. La Accuracy se utilizó como indicador global complementario.

\[
Recall = \frac{TP}{TP+FN}
\]

\[
F1 = 2 \cdot \frac{Precision \cdot Recall}{Precision + Recall}
\]

\[
Accuracy = \frac{TP+TN}{TP+TN+FP+FN}
\]

Donde $TP$ representa ataques minoritarios correctamente detectados, $FN$ representa ataques reales omitidos por el modelo, $FP$ representa registros no pertenecientes a la clase evaluada clasificados erróneamente como ataque y $TN$ representa registros correctamente descartados como no pertenecientes a la clase evaluada.

Las métricas computacionales fueron tiempo de entrenamiento, latencia de inferencia y consumo de memoria RAM. El tiempo de entrenamiento se midió en segundos, la latencia de inferencia en milisegundos y la memoria RAM en megabytes. Estas mediciones permitieron determinar si el modelo híbrido, además de mejorar la detección, resulta viable para escenarios de análisis de tráfico IoT.

\subsection{Prueba de hipótesis}

Para validar estadísticamente la diferencia entre el modelo híbrido y los modelos base se aplicó la prueba de rangos con signo de Wilcoxon. Esta prueba no paramétrica fue pertinente porque permite comparar mediciones relacionadas sin exigir normalidad en la distribución de las diferencias.

La hipótesis nula ($H_0$) estableció que no existen diferencias estadísticamente significativas entre el rendimiento del modelo híbrido y el rendimiento de los modelos base. La hipótesis alternativa ($H_1$) estableció que el modelo híbrido presenta diferencias significativas favorables en las métricas de detección de ataques minoritarios.

\[
H_0: M_{hibrido} = M_{base}
\]

\[
H_1: M_{hibrido} > M_{base}
\]

El criterio de decisión fue rechazar $H_0$ cuando el valor $p$ obtenido sea menor al nivel de significancia $\alpha = 0.05$. La prueba se aplicará sobre los resultados comparables de Recall, F1-Score y Accuracy, priorizando las clases minoritarias.

\subsection{Control de validez y replicabilidad}

Para proteger la validez interna, el rebalanceo se aplicó únicamente al conjunto de entrenamiento, evitando que las instancias sintéticas influyan en la evaluación final. Asimismo, se mantuvieron particiones independientes de entrenamiento, validación y prueba, y se utilizaron las mismas particiones para todos los modelos comparados.

Para fortalecer la validez externa, se utilizó un dataset estandarizado y reconocido en investigación de ciberseguridad IoT. Aunque los resultados dependen de las condiciones del CICIoT2023, el procedimiento puede replicarse sobre otros conjuntos de tráfico IoT con estructura similar.

Para asegurar replicabilidad, se documentaron las versiones de software, la estructura de datos, las particiones utilizadas, los criterios de inclusión y exclusión, los algoritmos evaluados, las métricas calculadas y el procedimiento de validación estadística. Además, se recomienda fijar semillas aleatorias en los algoritmos y registrar los hiperparámetros finales utilizados durante los experimentos.

\subsection{Resumen metodológico}

\begin{longtable}{p{4cm}p{8.5cm}}
\caption{Resumen del diseño metodológico de la investigación} \\
\toprule
\textbf{Elemento} & \textbf{Descripción} \\
\midrule
\endfirsthead
\toprule
\textbf{Elemento} & \textbf{Descripción} \\
\midrule
\endhead
Tipo de investigación & Aplicada. \\
Enfoque & Cuantitativo. \\
Nivel & Explicativo y comparativo. \\
Diseño & Experimental puro mediante simulación en entorno controlado. \\
Metodología de trabajo & CRISP-DM. \\
Población & Registros de tráfico benigno y malicioso del dataset CICIoT2023. \\
Muestra & Muestra censal del conjunto procesado: 7\,845\,673 registros y 47 variables. \\
Grupo experimental & Modelo híbrido Stacking con XGBoost, LightGBM, entropía de Shannon y Árbol de Decisión. \\
Grupo de control & Random Forest y XGBoost sin arquitectura de apilamiento. \\
Técnica de balanceo & Borderline-SMOTE aplicado únicamente al conjunto de entrenamiento. \\
Métricas predictivas & Recall, F1-Score y Accuracy. \\
Métricas computacionales & Tiempo de entrenamiento, latencia de inferencia y consumo de memoria RAM. \\
Prueba estadística & Rangos con signo de Wilcoxon, con $\alpha = 0.05$. \\
\bottomrule
\end{longtable}

\end{document}
